\documentclass[11pt]{article}

% --- Encoding and fonts (pdfLaTeX-friendly; compiles on Overleaf out of the box) ---
\usepackage[T1]{fontenc}
\usepackage[utf8]{inputenc}
\usepackage{lmodern}

% --- Layout ---
\usepackage[margin=1in]{geometry}
\usepackage{microtype}
\usepackage{parskip}
\usepackage{enumitem}
\setlist[itemize]{leftmargin=1.4em, itemsep=0.2em, topsep=0.3em}
\usepackage{titlesec}
\titleformat{\section}{\large\bfseries}{\thesection}{0.6em}{}
\titleformat{\subsection}{\normalsize\bfseries}{\thesubsection}{0.6em}{}

\usepackage{xcolor}
\usepackage[hidelinks]{hyperref}
\hypersetup{
  pdftitle = {Never Lose Your Place: Understanding as a Shared Build Artifact},
  pdfauthor = {Adam Carlson}
}

\newcommand{\vs}{\textit{Vibe Scribe}}

\title{\textbf{Never Lose Your Place}\\[0.45em]
  {\large Understanding as a Shared Build Artifact:\\[0.15em]
  A Paradigm for AI-Assisted Development}}
\author{Adam Carlson}
\date{June 2026}

\begin{document}
\maketitle

\begin{abstract}
\noindent
AI-assisted development has quietly inverted a long-standing assumption: that a developer
understands a system before building it, or can reconstruct that understanding by reading the
code afterward. When a machine writes code faster than any person can internalize it, both halves
of that assumption break. But the gap it opens is not the human's alone --- the assistant, too,
loses its grasp of a system across context-window limits, conversational compaction, and session
boundaries. This paper names the resulting condition, the understanding gap suffered by human and
AI in complementary ways, and proposes a paradigm for closing it: treat understanding as a
\emph{shared} build artifact, manufactured continuously and automatically as a side effect of the
work, into a persistent, plain-text substrate that both the human and the AI read, that keeps the
two aligned to a single model of the system, and that is kept honest against the source of truth.
We state the paradigm's commitments, present an open-source existence proof (\vs{}), mark its
deliberate exclusions and honest limitations, sketch how it should aid debugging and iteration,
and offer predictions by which the paradigm can be judged. We make no efficacy claims: this is a
paradigm and a working demonstration, not a measured result.
\end{abstract}

\section{A place lost}

On an ordinary afternoon, a development machine restarted itself in the middle of a working
session. The change in progress was delicate and half-finished, and a dozen small decisions behind
it lived only in the back of the developer's mind and in the volatile context of an AI assistant.
By every normal expectation, that thread was gone.

It was not. A single plain-text file --- a running journal the assistant had been keeping as a
side effect of the work --- recorded what had changed, why, and what remained. The next session
read that file and resumed at the exact point of interruption, as if nothing had happened.

Note what was recovered. Not only the developer's place in the work, but the \emph{assistant's}
model of it --- the architecture, the decisions, the dead ends already ruled out. Both had lost
the thread; one plain-text file returned it to both. That small, two-sided recovery is the seed of
this paper, because the relief it produced points at a problem most developers now feel but few
have named.

\section{The phenomenon nobody has named}

Call it the \emph{understanding gap}: the distance between the code that exists in a project and
the understanding of that code that exists in anyone's head.

The gap has two axes. The first is \textbf{continuity} --- you understood the system a moment ago,
and a boundary severed the thread. The second is \textbf{comprehension} --- you may never have
held the thread at all, because the system was assembled faster than you could form a model of it.
Neither is a deficiency of skill: a newcomer meets the gap as a wall, an expert as drag, the
accumulating weight of code they directed but never fully absorbed.

And the gap is not the human's alone. A human forgets slowly and selectively --- across time,
divided attention, staff turnover, or never having grasped the system in the first place. An AI
assistant forgets abruptly and completely: the context window fills, the conversation is compacted
into a lossy summary, the session ends, and its model of the architecture, of what was tried, of
the progress made, is simply gone at the next turn. The two forget in \emph{complementary} ways,
and the boundaries that sever the thread divide along the same line. Time, handoff, and inattention
are where the human loses the thread; the context window, compaction, and the session boundary are
where the AI loses its; a reboot or a tool switch takes both at once. Each is a place someone loses
their place.

\section{The old paradigm, and why it is breaking}

Two assumptions sit, unstated, beneath conventional practice. The first is that \emph{understanding
precedes code}: you grasp a design, then you implement it. The second is that understanding
\emph{can be reconstructed from} code: when the grasp fades, you read the source and recover it.
Documentation, under this paradigm, is a separate and largely human effort --- written after the
fact, for other humans, and left to decay.

AI-assisted development breaks both assumptions at once. When a model generates implementation
faster than a person can internalize it, ``grasp, then build'' collapses: the building outruns the
grasping. And ``read it back'' does not scale to recover what was never deliberately recorded ---
intent, rejected alternatives, the reason a thing is shaped the way it is. Code is an excellent
record of \emph{what}; it is silent on \emph{why}. The old paradigm assumed a human pace of
authorship, and a single mind to hold the model. AI-assisted work violates both: the model of the
system is now split across two collaborators, neither of whom retains it durably.

\section{The new paradigm: understanding as a shared build artifact}

We propose a different arrangement. In AI-assisted development, understanding should be
\textbf{manufactured continuously, as a shared build artifact} --- produced alongside the code, by
the same machine that writes the code, and deposited into a persistent substrate that both the
human and the AI read, kept true automatically as the work proceeds.

Two things follow from the word \emph{shared}. The first is the inversion: understanding stops
being a \emph{prerequisite} the developer is assumed to bring, and stops being an \emph{archaeology
project} to be undertaken later. It becomes an \emph{output} --- a first-class artifact of the
build, accruing continuously and available on demand.

The second is less obvious and more powerful. Because both parties read the \emph{same} artifact,
they are aligned to the \emph{same} model of the system. This is more than two separately-informed
parties; it is a single source of truth for a shared mental model. It converts shared understanding
from something renegotiated every session --- you explain, the assistant confirms, you both proceed,
and tomorrow you do it again --- into something durable and co-owned: established once and maintained
continuously. The relationship shifts from \emph{a user instructing a tool} to \emph{two
collaborators reasoning from a common, living model.}

``Never lose your place'' is the felt promise of this arrangement --- and it is a promise to both
parties at once. One honest qualifier: the artifact is a shared external \emph{reference} that keeps
the two aligned; it is not identical internal cognition. The record holds the common ground; a
briefing is where the two reconcile against it --- and the paradigm can make that reconciliation
\emph{active}, prompting it at the moments the system has just changed, so alignment is restored
exactly when it is most at risk of slipping.

\section{What the paradigm commits you to}

A paradigm is more than a slogan; it implies constraints. Several follow directly from the claim
above, independent of any particular implementation.

\begin{itemize}
  \item \textbf{A shared substrate.} The record is for the human \emph{and} the machine --- plain
  text in the project's own repository, readable by a person, parseable by any tool, free of
  lock-in. Because it is one substrate read by both, it doubles as common ground: the two parties
  reason from the same names, the same map, the same recorded decisions.

  \item \textbf{Two shapes of memory.} Continuity and comprehension want opposite structures. A
  \emph{map} --- revised in place --- describes the system as it is now; this is what lets you grasp
  it. A \emph{log} --- append-only, never edited --- preserves the temporal thread: not every edit,
  but the \emph{decisions}, the \emph{reasons} behind them, the \emph{alternatives rejected}, and the
  \emph{surprises}. It records why the system came to be the way it is --- the one thing the code, a
  record only of its latest state, can never recover --- and it is written for a reader who was not
  present, pre-eminently the next AI session. Its most valuable service is to keep that reader from
  \emph{repeating itself}: re-exploring a dead end, re-proposing a rejected approach, re-introducing a
  fixed bug. Same maintenance habit, opposite update semantics; conflating them destroys both.

  \item \textbf{Maintenance as a side effect.} If keeping the record current is a separate task the
  human must remember, it will not happen. The substrate must be maintained automatically, as a
  byproduct of the work itself.

  \item \textbf{Honesty about staleness.} A record that cannot distinguish what it has verified from
  what may have drifted is worse than none, because it invites false confidence. Freshness must be a
  first-class signal the record carries and surfaces.

  \item \textbf{Self-verification.} The understanding must be checkable against the source of truth.
  The machine proposes the record; the code verifies it --- the record audits itself against the code
  rather than decaying unchecked.

  \item \textbf{A shared vocabulary.} A person cannot navigate a system they have no names for. The
  substrate coins memorable, human-readable names for the system's parts, anchored to real code, so a
  single handle indexes across the map, the history, and the source --- and means the same thing to
  the human and the AI.
\end{itemize}

\section{An existence proof: \vs{}}

\vs{} is one realization of the paradigm, open-source and runnable today.\footnote{Source and
documentation: \url{https://github.com/adamccarlson/vibe-scribe} (MIT-licensed).} It is not the
paradigm; it is evidence the paradigm is buildable rather than aspirational.

It keeps two plain-markdown documents in the repository: a \emph{dev journal} (the log) and a
\emph{system overview} (the map, which carries a \emph{Lexicon} of coined names anchored to the
code). The documents are maintained automatically: a lightweight, fail-open hook notices when source
code changed but the documents did not, and quietly asks the assistant to record the change before
finishing. The map carries per-entry verification dates, so stale sections announce themselves; an
audit command re-checks the map against the code and refreshes or flags it. A briefing command reads
the documents back in plain language --- with small text diagrams --- so the project can explain
itself to its author, and so the assistant can re-ground itself at the start of a session and
reconcile with the user on the current state. After a significant change, the assistant can go
further and \emph{offer} such a briefing unprompted, framed as an intent check --- ``does this match
what you had in mind?'' --- surfacing any gap between what was intended and what was built while it is
fresh and cheap to correct. The same core runs across three AI coding tools
(Claude Code, OpenAI Codex, and Cursor) from one shared implementation, because the artifact it
produces is tool-agnostic text. Each of these is a tenet from the previous section, made concrete;
none requires research-grade machinery, and the result belongs to the project, not to any vendor.

\section{What the paradigm excludes}

A lens is defined as much by what it leaves out. \vs{} declined a formal knowledge graph of the
system's components: the relationships it would encode are already captured, in human-readable form,
by the map and its vocabulary, and a machine-only graph is a lossy second copy of something the code
already encodes precisely and that static analysis can extract for free. A graph earns its keep when
the graph \emph{is} the product (as in a graph of ideas with no other representation), not when it
shadows source that already exists. For the same reason it declined a rigid relationship ontology,
preferring open, plain-language relationship labels, and declined rendered or binary diagrams in
favor of plain text that diffs, travels, and reads anywhere. The paradigm favors the representation a
human can read and a tool can verify over the one that is merely machine-precise.

\section{Honest positioning}

This paradigm stands on familiar ground, and says so --- but the prior art divides almost exactly
along the line this paper has drawn. Engineering notebooks, architecture decision records, runbooks,
README-driven development, living documentation, and domain-driven design's notion of a shared
ubiquitous language all serve \emph{human} understanding. The ``memory'' files, context-management
schemes, and retrieval-over-conversation emerging around coding agents serve \emph{machine}
understanding. Each side addresses one party's forgetting and ignores the other's.

What is new here is not a component but a \emph{unification}: a single substrate that serves both,
maintained by the AI, read by both, and --- the decisive part --- that keeps the two \emph{aligned}
to the same model rather than letting their understandings drift apart. The recurring failure of the
prior art is not conception but sustainability: the records rot because maintenance is manual and
staleness is invisible. The contribution is to move that discipline into the build loop, to make the
record honest about its own freshness, and to make it the common ground of a human--AI pair.

We claim no measured benefit. \vs{} is new; it has not been evaluated; it leans in part on the
assistant following instructions; verification is harder to sustain as a system grows; and its
automatic enforcement has not been exercised live in every supported tool. This paper presents a
paradigm and a working demonstration of it --- not evidence that it improves any outcome. That
honesty is itself a commitment of the paradigm: a record, including this one, should not claim more
confidence than it has earned.

\section{What it should make easier: debugging and iteration}

A paradigm should pay rent in practice, and this one makes specific --- though, we stress, as-yet-%
unmeasured --- promises about the two activities that consume most of a developer's time: fixing what
is broken and improving what works. At least five mechanisms are distinct.

\begin{itemize}
  \item \textbf{Localization.} Most defects arrive with recent changes, and the log is precisely a
  record of recent change and its rationale --- the first question in debugging, answered directly.

  \item \textbf{Blast-radius reasoning.} The map's relationships --- which component owns which data,
  what calls what --- bound where a symptom can originate and reveal what a fix might break.

  \item \textbf{Compounding failure memory.} A diagnosed fault, once recorded, is thereafter looked
  up rather than re-diagnosed, so debugging knowledge accumulates instead of evaporating.

  \item \textbf{Non-repetition.} Most pointed for the AI: because the record holds the \emph{why} and
  the rejected paths, neither party ``fixes'' a problem by reverting to an approach already known to
  fail, or by violating an invariant an earlier decision was protecting.

  \item \textbf{Architecture-respecting change.} An accurate, shared model lets improvements fit the
  existing design rather than fight it, and lets human and AI reason about the change from the
  \emph{same} model in the \emph{same} vocabulary --- fewer mismatches of the form ``I thought that
  component did $X$.''
\end{itemize}

These are affordances the paradigm should provide, not results it has demonstrated. Whether they hold
is what the experience report below would test.

\section{What the paradigm predicts}

A paradigm worth the name makes claims that could turn out false. If understanding-as-a-shared-build-%
artifact is the right frame, several things should follow. The binding constraint in AI-assisted
development should shift from \emph{writing} code to \emph{comprehending and handing off} code.
Tooling that manufactures understanding should become as ordinary as version control. AI coding
agents should come to depend on durable external substrates that they themselves maintain --- and the
artifact that serves the human's comprehension and the one that serves the agent's continuity should
converge into a single shared record, rather than living in two separate tools. And teams that skip
all this should accumulate a compounding ``comprehension debt'': shippable systems that nobody, human
or machine, can safely change.

The paradigm is also disconfirmable. If models become capable enough that code never needs to be
understood by a person again, or if static analysis can reconstruct intent --- not just structure ---
fully and on demand, then manufacturing understanding is wasted effort and the paradigm is moot. The
honest next step is neither argument nor assertion but use: run the approach on real projects over
real time, and report what actually happens. That experience report, not this paper, is where evidence
begins.

\section{The smallest promise, the largest stakes}

A machine rebooted, and neither the developer nor the assistant lost their place. That is the smallest
possible version of the claim, and it is also the whole of it. As AI writes more of the code,
understanding stops arriving for free, and it stops being something either mind reliably holds alone.
It has to be manufactured --- continuously, automatically, honestly, and \emph{together} --- or it
simply will not be there when the thread breaks, as it always eventually does.

Here is one paradigm for manufacturing it, and one proof that it can be built. Never lose your place
--- for either of you.

\section*{Acknowledgments}

\vs{} was designed and implemented, and this paper drafted, in close collaboration with an AI coding
assistant --- which is, fittingly, the very condition the paradigm addresses.

\end{document}
